\chapter*{Angabe zu verwendeten KI-basierter Elektronischer Hilfsmittel}
\vspace{-6em}{\Huge\color{dkfzgray2} Information on the use of AI-based tools}
\addcontentsline{toc}{chapter}{Angabe zu verwendeten KI-basierter Elektronischer Hilfsmittel}

Zur Dokumentation der verwendeten Hilfsmittel ist der schriftlichen Ausarbeitung ein besonderer Anhang hinzugefügt, der eine Liste und Beschreibung aller verwendeter KI-basierter Hilfsmittel enthält. Der besondere Anhang zur Dokumentation der verwendeten Hilfsmittel erfüllt folgende Kriterien:

\emph{\small\color{dkfzgray2} To document the tools used, a special appendix has been added to the written report, which contains a list and description of all AI-based tools used. The special appendix to document the tools used meets the following criteria:}

\begin{enumerate}
    \item Auflistung der Ziele, für die die KI-basierten Hilfsmittel in der vorliegenden Arbeit eingesetzt wurden.\\
    \emph{\small\color{dkfzgray2} List of the goals for which the AI-based tools were used in this work.}
    \item Dokumentation der Verwendungsweise der KI-basierten Hilfsmittel.\\
    \emph{\small\color{dkfzgray2} Documentation of how AI-based tools are used.}
    \item Nennung der Kapitel und Abschnitte der vorliegenden Arbeit, in denen die KI-basierten Hilfsmittel eingesetzt wurden, um Inhalte zu erzeugen.\\
    \emph{\small\color{dkfzgray2} Identification of the chapters and sections of this work in which AI-based tools were used to generate content.}
\end{enumerate}

Der Gebrauch dieser Hilfsmittel inklusive Art, Ziel und Umfang des Gebrauchs wurde mit meinem offiziellen Betreuer / meiner offiziellen Betreuerin \textbf{TODO NAME OF SUPERVISOR} abgesprochen.

\emph{\small\color{dkfzgray2} The use of these aids, including the type, purpose and extent of use, has been agreed with my official supervisor \textbf{TODO NAME OF SUPERVISOR}.}
\newpage
Mir ist bewusst, dass insbesondere der Versuch einer nicht dokumentierten Nutzung KI- basierter Hilfsmittel als Täuschungsversuch zu werten ist:

\emph{\small\color{dkfzgray2} I am aware that in particular the attempt to use AI-based tools without documentation is to be regarded as an attempt at deception:}

Gem. § 16 Abs. 2 der Promotionsordnung „Dr. med. / dent.“:
„Ergibt sich vor Aushändigung der Promotionsurkunde, dass der Kandidat/die Kandidatin bei einer Promotionsleistung getäuscht hat, so können einzelne oder alle Promotionsleistungen für ungültig erklärt werden. In schweren Fällen kann die Zulassung zum Promotionsverfahren zurückgenommen werden.“

\emph{\small\color{dkfzgray2} According to § 16 paragraph 2 of the doctoral regulations for "Dr. med. / dent.":
"If it emerges before the doctoral certificate is issued that the candidate has cheated in a doctoral achievement, individual or all doctoral achievements can be declared invalid. In serious cases, admission to the doctoral process can be withdrawn.}

Und § 16 Abs. 2 der Promotionsordnung „Dr. sc. hum.“:
„Ergibt sich vor Aushändigung der Promotionsurkunde, dass der Doktorand / Doktorandin bei einer Promotionsleistung getäuscht hat, so kann der Promotionsausschuss diese Promotionsleistung oder alle bisher erbrachten Promotionsleistungen für ungültig erklären. In besonders schweren Fällen kann der Promotionsausschuss die Annahme als Doktorand / Doktorandin endgültig widerrufen.“

\emph{\small\color{dkfzgray2} And § 16 paragraph 2 of the doctoral regulations “Dr. sc. hum.”:
“If it emerges before the doctoral certificate is issued that the doctoral candidate has cheated in a doctoral achievement, the doctoral committee can declare this doctoral achievement or all doctoral achievements to date invalid. In particularly serious cases, the doctoral committee can permanently revoke the acceptance as a doctoral candidate.”}

\begin{tabular}{p{3cm}p{5cm}p{4cm}}
      &  & \\
      & & \\
      &Heidelberg, \today & \hrulefill\\
      & & \textit{TODO: name of doctoral student}
\end{tabular}
